\chapter{Introduction}
\label{chap:introduction}

\section{Motivation}
\label{sec:intro-motivation}

Jupyter notebooks combine code, explanatory text, and results in a single
document. This makes them a popular way to share computational research,
especially in data-heavy fields such as biomedicine. A notebook published
alongside a paper is meant to let another researcher see, and re-run, exactly
how a result was produced.

In practice, many published notebooks do not run at all when someone else
tries them. Pimentel et al.~\cite{pimentel2019msr} collected 1.4~million
notebooks from GitHub and attempted to execute a large subset of them: only
24.11\% ran without errors, and only 4.03\% reproduced the results reported by
their original authors. Samuel and Mietchen~\cite{samuel2024gigascience}
found a similar pattern specifically for notebooks accompanying biomedical
publications on PubMed Central, the domain this thesis also builds on. These
numbers show that non-reproducible notebooks are not a rare exception; they
are the common case.

A large share of these failures is caused by dependencies, not by mistakes in
the underlying research logic. A notebook is written against a particular set
of installed Python packages. When it is later run in a different
environment, a package may be missing entirely, or a newer version of a
package may no longer provide the exact function the notebook expects. Both
situations produce an error before the notebook can finish, even though the
scientific code itself may be correct. The resulting error messages are often
short and technical, for example \texttt{ModuleNotFoundError: No module named
'sklearn'}, which is unhelpful to a researcher who is not a software
developer and has no way to know that the fix is to install a package called
\texttt{scikit-learn} rather than \texttt{sklearn}.

\section{The FAIR Jupyter Pipeline and Its Remaining Gap}
\label{sec:intro-gap}

The FAIR Jupyter project has studied this problem at scale. Samuel and
Mietchen~\cite{samuel2024gigascience} built a pipeline that downloads
notebooks referenced by biomedical publications, executes them, and records
whether each notebook ran and, if not, what error occurred. The results were
organised into a queryable knowledge graph~\cite{samuel2024fairkg} so that
patterns in notebook reproducibility could be studied directly. A more recent
extension of the same pipeline~\cite{samuel2025docker} runs every repository
inside its own Docker container before execution. Containerisation removes
many failures that were really about a missing system tool or an
incompatible operating system, rather than about the research code.

Even this containerised pipeline stops at the point of \emph{detecting} a
failure. When a notebook still does not run inside its container, the
pipeline records the exception type and message and moves on. It does not
explain the failure to a human reader, and it does not attempt a fix. Every
notebook whose only problem is a missing or outdated Python package is left
exactly as broken as it was found, even though such problems are, in
principle, straightforward to repair by installing or pinning the right
package version.

This is the gap this thesis addresses: turning a detected, dependency-related
notebook failure into an explained and, where possible, repaired one,
automatically and without a person in the loop for every single notebook.

\section{Contribution of This Thesis}
\label{sec:intro-contribution}

This thesis adds a repair layer on top of the existing Docker-based FAIR
Jupyter pipeline. The pipeline's own contribution ends at logging a failure;
this thesis begins exactly there. Given a notebook execution that failed with
a dependency-related error, the added layer performs five steps:

\begin{enumerate}
  \item \textbf{Classify} the failure: is it really a dependency problem, and
        if so, is it a missing package or an incompatible version?
  \item \textbf{Explain} the failure in plain language, so that a researcher
        without a programming background can understand what went wrong.
  \item \textbf{Propose a repair}, using an open, locally-run large language
        model (LLM) whose suggestion is grounded in live package metadata
        from the Python Package Index (PyPI) rather than left to guess a
        version number from training data.
  \item \textbf{Apply the repair} inside an environment equivalent to the one
        the notebook originally failed in, and re-run the whole notebook from
        the beginning.
  \item \textbf{Record the outcome}: whether the notebook now runs further,
        still fails, or could not be meaningfully tested.
  \end{enumerate}

Grounding the repair suggestion in live PyPI data follows PLLM, the closest
prior system to this one. Bartlett et
al.~\cite{bartlett2025pllm} show that retrieval-augmented generation (RAG)
against live package data measurably improves an LLM's dependency-repair
accuracy compared to letting the model rely on what it happens to remember
from training. This thesis adopts the same principle and applies it inside a
reproducibility pipeline rather than to standalone Python scripts, and adds
an explanation step and a re-execution-based validation step that PLLM does
not have.

\subsection{Objectives}
\label{sec:intro-objectives}

The work is organised around six objectives, first defined in the thesis
vision document and used throughout this text to relate individual
components back to the overall goal.

\begin{table}[h]
  \centering
  \begin{tabular}{|l|p{10.5cm}|}
    \hline
    \textbf{Objective} & \textbf{Description} \\
    \hline
    O1 & Generate a plain-language explanation of a dependency-related failure. \\
    \hline
    O2 & Generate an automated, PyPI-grounded fix suggestion using an open LLM. \\
    \hline
    O3 & Apply the suggested fix and validate it by re-executing the notebook. \\
    \hline
    O4 & Compile a structured, reusable benchmark dataset of failures, fixes, and outcomes. \\
    \hline
    O5 & Integrate the above into the existing FAIR Jupyter pipeline without manual intervention. \\
    \hline
    O6 & Record repair outcomes as RDF triples in the FAIR Jupyter knowledge graph. \\
    \hline
  \end{tabular}
  \caption{Thesis objectives, as defined in the vision document.}
  \label{tab:intro-objectives}
\end{table}

\subsection{Status of This Thesis}
\label{sec:intro-status}

This thesis reports work completed through the result-logging and
knowledge-graph implementation stage of a staged
design-implementation-evaluation plan; \Cref{chap:methodology} explains,
component by component,
how each stage was actually built. As of this stage:

\begin{itemize}
  \item Objectives O1, O2, and O3 are \textbf{implemented and evaluated}: a
        classification and explanation component, a PyPI-grounded repair
        proposal component, and a fix-application-and-re-execution component
        all exist as working, unit-tested code, and have since been run
        together, unattended, over the full 187-row evaluation split
        described in \Cref{chap:problem-dataset}, using a real Ollama model,
        the real PyPI registry, and a real Docker daemon throughout. No
        notebook in this split reached completely successful re-execution
        within the two-round repair budget, but 84.9\% of the 73 real
        repair attempts made across both rounds removed the specific
        dependency error they targeted, and every language-model repair
        response obtained during the run was schema-valid and PyPI-grounded.
        \Cref{chap:validation} reports this evaluation in full.
  \item Objectives O4 and O6 -- the reusable benchmark dataset and
        knowledge-graph enrichment -- are \textbf{implemented and validated
        at pilot scale}: a result-logging component persists one row per
        repair attempt to a SQLite table, and an RDF mapping links each
        logged attempt into the existing FAIR Jupyter knowledge graph. The
        benchmark table has since been populated at the scale of the full
        evaluation run above; enriching the knowledge graph from that larger
        run, rather than from the pilot sample checked in
        \Cref{sec:v-i6}, remains future work.
  \item Objective O5, full pipeline integration, is \textbf{implemented and
        evaluated at full scale}: a single program
        (\texttt{scripts/run\_pipeline.py}) runs all five repair-layer
        components, plus a bounded, max-two-round repair loop, over a given
        slice of the dataset automatically, sharing one run identifier
        across every stage it calls. Beyond the small real pilot that first
        validated this orchestrator (\Cref{sec:v-i7}), it has now driven the
        complete 187-row evaluation run referenced above.
        \Cref{chap:future-work} states exactly what remains.
  \item The evaluation reported in \Cref{chap:validation} is the
        thesis's final evaluation, not only a pilot: it covers the complete
        187-row evaluation split, once, under the same configuration
        throughout. It measures an overall repair success rate together
        with the diagnostic metrics needed to interpret it -- the rate at
        which individual targeted dependency errors were resolved, the rate
        at which the repair agent could reach the language model at all
        given its current package-mapping coverage, and the rate of
        infrastructure faults unrelated to repair quality. \Cref{chap:validation}
        reports each of these in turn, together with what they do and do
        not show about the system.
\end{itemize}

This distinction between what has been implemented and validated, and what
is still planned, is kept explicit throughout this document rather than
blurred for the sake of a stronger-sounding narrative.

\section{Thesis Structure}
\label{sec:intro-structure}

\Cref{chap:related-work} reviews prior work on notebook reproducibility,
automated dependency resolution, and LLM-based program repair, and states
the specific gap this thesis fills. \Cref{chap:problem-dataset} describes the
scale of the reproducibility problem and how the concrete working dataset for
this thesis was obtained. \Cref{chap:architecture} presents the system
architecture: the five components introduced above, their contracts, and the
data they exchange. \Cref{chap:methodology} is the core of the thesis: it
explains, component by component, how each part was actually built, the
decisions made along the way, and where the implementation ended up
differing from the original design. \Cref{chap:validation} reports the
automated test suite, the real-system pilot results obtained during
implementation, and the final evaluation over the full 187-row evaluation
split. \Cref{chap:future-work} separates what is finished from what remains.
\Cref{chap:conclusion} closes the thesis with a summary of the contribution.
